% !TeX spellcheck = pt_BR

\chapter*{Introdução ao ato final}
\addcontentsline{toc}{chapter}{Introdução ao ato final}

 
Ao longo das partes anteriores, desenvolvemos duas linguagens
distintas para descrever e compreender o movimento dos corpos.
A primeira, construída na Parte~I, foi a Cinemática: aprendemos a
\emph{descrever} movimentos, tanto na linguagem escalar ---
adequada para movimentos em uma dimensão --- quanto na linguagem
vetorial, necessária para tratar movimentos em duas ou três
dimensões. A segunda linguagem foi construída na Parte~II, a
Dinâmica: primeiro na sua formulação vetorial, baseada nas Leis
de Newton e na noção de força como causa da variação do movimento;
depois na sua formulação escalar, baseada nos conceitos de trabalho
e energia, capaz de responder perguntas sobre os estados inicial e
final de um sistema sem que precisemos conhecer todos os detalhes
do percurso.
 
Essas duas formulações da dinâmica não são concorrentes: são
complementares. Como já discutimos, cada uma ilumina aspectos que
a outra deixa na sombra. A formulação newtoniana é insubstituível
quando precisamos conhecer as forças instantâneas, as trajetórias e
o que ocorre ao longo do caminho. A formulação energética é
imbatível quando queremos comparar estados extremos de um sistema
sem nos preocupar com o que aconteceu entre eles.
 
% ---------------------------------------------------------------
%  TEXTO DO PARALELO ESCALAR / VETORIAL
%  (integrado na abertura da Parte III)
% ---------------------------------------------------------------
 
Vale, antes de prosseguir, dar um passo atrás e contemplar a
arquitetura que construímos ao longo deste volume --- porque ela
carrega, por si só, uma lição sobre a estrutura da mecânica
clássica.
 
A Parte~I, a Cinemática, se divide em dois blocos que o leitor
já conhece bem: a \textbf{Cinemática Escalar} (Capítulo~1) e a
\textbf{Cinemática Vetorial} (Capítulos~2 e~3). Começamos pelo
escalar --- movimentos em uma dimensão, tratados algebricamente
--- porque ele é mais simples, e porque o caráter vetorial de
grandezas como velocidade e aceleração pode, nesse contexto, ser
substituído por uma simples atribuição de sinais. Só quando
precisamos descrever movimentos que \emph{viram} --- trajetórias
curvas, movimentos em duas ou três dimensões --- somos forçados a
recorrer aos vetores.
 
A Parte~II, a Dinâmica, repete a dualidade, mas em ordem
invertida: começamos pela \textbf{Dinâmica Vetorial} (Capítulo~4)
--- as Leis de Newton, forças, acelerações, momento linear --- e
só então desenvolvemos a \textbf{Dinâmica Escalar} (Capítulo~5)
--- trabalho, energia, potência, conservação. A ordem se inverte
porque a lógica agora é outra: para entender o que \emph{causa} o
movimento, precisamos das forças --- e forças são vetores. A
formulação energética, que dispensa as direções e os percursos,
só se torna inteligível depois que a formulação vetorial está bem
estabelecida, porque é dela que emerge.
 

 
 
% ---------------------------------------------------------------
%  RETOMADA E ANÚNCIO DOS CAPÍTULOS
% ---------------------------------------------------------------
 
Chegamos agora à Parte~III munidos dessas linguagens. Poderíamos
chamá-la de \emph{síntese}, mas o nome escolhido ---
\emph{Aplicações} --- carrega um sentido mais preciso. Não se trata
de tópicos menores, apêndices da teoria ``real''. Os problemas
desta parte são, em vários aspectos, os mais ricos e exigentes de
todo o volume. Chamamo-los de aplicações no sentido em que o
próprio Newton usou a palavra ao abrir o terceiro livro dos
\emph{Principia}: ``passo agora a demonstrar o sistema do
mundo''.\footnote{Newton, I. \textit{Philosophiæ Naturalis
Principia Mathematica}, Livro~III (\textit{De Mundi Systemate}),
1687.} Assim como Newton reservou seus dois primeiros livros para
construir a teoria e o terceiro para aplicá-la aos céus,
reservamos nossas duas primeiras partes para construir as linguagens
da mecânica e esta para aplicá-las a dois domínios que, juntos,
cobrem uma extensão extraordinária da realidade física.
 
O primeiro é a \textbf{gravitação} --- o problema que, literalmente,
motivou Newton a escrever os \emph{Principia}. A gravitação é o
terreno em que as duas linguagens se mostram mais complementares:
a vetorial é indispensável para tratar forças gravitacionais e
acelerações centrípetas; a escalar é a ferramenta natural para
tratar energia orbital e velocidade de escape. Quem só sabe usar
uma das duas vai conseguir resolver metade dos problemas.
 
O segundo é a \textbf{dinâmica dos fluidos} --- o estudo do
comportamento de líquidos e gases, em repouso e em movimento.
Aqui as ferramentas conceituais que desenvolvemos serão postas à
prova de uma forma nova: os fluidos não são corpos pontuais nem
rígidos, e o tratamento adequado exige que saibamos transitar entre
pressão (grandeza escalar, distribuída por área) e força (grandeza
vetorial, localizada).
 
O terceiro ato é, portanto, o momento em que a mecânica deixa de ser
um conjunto de técnicas e se torna uma forma de ver o mundo.